% prezentace.tex -- Lekce 10: Volný projekt
% preamble-prez.tex -- sdílená preamble pro prezentace (beamer)
% Includuje se z ucitel/lekceXX/prezentace.tex jako:
%   % preamble-prez.tex -- sdílená preamble pro prezentace (beamer)
% Includuje se z ucitel/lekceXX/prezentace.tex jako:
%   % preamble-prez.tex -- sdílená preamble pro prezentace (beamer)
% Includuje se z ucitel/lekceXX/prezentace.tex jako:
%   \input{../spolecne/preamble-prez.tex}

\documentclass[aspectratio=169]{beamer}

% --- Jazyk a font ---
\usepackage{polyglossia}
\setmainlanguage{czech}
\usepackage{fontspec}

% --- Téma ---
\usetheme{default}
\usecolortheme{default}

\definecolor{microbit-blue}{HTML}{1E4D8C}
\definecolor{microbit-lightblue}{HTML}{D6E4F7}
\definecolor{code-bg}{HTML}{F0F0F0}

\setbeamercolor{frametitle}{bg=microbit-blue, fg=white}
\setbeamercolor{title}{fg=microbit-blue}
\setbeamercolor{block title}{bg=microbit-blue, fg=white}
\setbeamercolor{block body}{bg=microbit-lightblue}
\setbeamercolor{itemize item}{fg=microbit-blue}
\setbeamercolor{itemize subitem}{fg=microbit-blue!70}
\setbeamerfont{frametitle}{size=\large, series=\bfseries}
\setbeamertemplate{navigation symbols}{}
\setbeamertemplate{footline}{%
  \leavevmode%
  \hbox{%
    \begin{beamercolorbox}[wd=.5\paperwidth,ht=2.5ex,dp=1.125ex,leftskip=6pt]{author in head/foot}%
      \usebeamerfont{author in head/foot}\textcolor{gray}{\small Úvod do Pythonu na Micro:bitu}
    \end{beamercolorbox}%
    \begin{beamercolorbox}[wd=.5\paperwidth,ht=2.5ex,dp=1.125ex,rightskip=6pt,right]{date in head/foot}%
      \usebeamerfont{date in head/foot}\textcolor{gray}{\small\insertframenumber\,/\,\inserttotalframenumber}
    \end{beamercolorbox}}%
}

% --- Česky korektní uvozovky ---
\usepackage{csquotes}
\newcommand{\uv}[1]{\enquote{#1}}

% --- Zdrojový kód (Python) ---
\usepackage{listings}
\lstset{
  language=Python,
  basicstyle=\ttfamily\small,
  backgroundcolor=\color{code-bg},
  frame=single,
  framerule=0pt,
  xleftmargin=6pt,
  xrightmargin=6pt,
  breaklines=true,
  showstringspaces=false,
  keywordstyle=\color{microbit-blue}\bfseries,
  stringstyle=\color{teal},
  commentstyle=\color{gray}\itshape,
  literate=
    {á}{{\'a}}1 {é}{{\'e}}1 {í}{{\'i}}1 {ó}{{\'o}}1 {ú}{{\'u}}1
    {ů}{{\r{u}}}1 {č}{{\v{c}}}1 {š}{{\v{s}}}1 {ž}{{\v{z}}}1
    {Á}{{\'A}}1 {É}{{\'E}}1 {Í}{{\'I}}1 {Ó}{{\'O}}1 {Ú}{{\'U}}1
    {Č}{{\v{C}}}1 {Š}{{\v{S}}}1 {Ž}{{\v{Z}}}1 {ň}{{\v{n}}}1 {ř}{{\v{r}}}1,
}

% --- Zvýraznění pojmů ---
\newcommand{\pojem}[1]{\textbf{\textcolor{microbit-blue}{#1}}}
\newcommand{\pyinline}[1]{\texttt{#1}}

% --- Grafika a diagramy ---
\usepackage{tikz}

% --- Ostatní ---
\usepackage{graphicx}
\usepackage{hyperref}
\hypersetup{colorlinks=true, urlcolor=microbit-blue}


\documentclass[aspectratio=169]{beamer}

% --- Jazyk a font ---
\usepackage{polyglossia}
\setmainlanguage{czech}
\usepackage{fontspec}

% --- Téma ---
\usetheme{default}
\usecolortheme{default}

\definecolor{microbit-blue}{HTML}{1E4D8C}
\definecolor{microbit-lightblue}{HTML}{D6E4F7}
\definecolor{code-bg}{HTML}{F0F0F0}

\setbeamercolor{frametitle}{bg=microbit-blue, fg=white}
\setbeamercolor{title}{fg=microbit-blue}
\setbeamercolor{block title}{bg=microbit-blue, fg=white}
\setbeamercolor{block body}{bg=microbit-lightblue}
\setbeamercolor{itemize item}{fg=microbit-blue}
\setbeamercolor{itemize subitem}{fg=microbit-blue!70}
\setbeamerfont{frametitle}{size=\large, series=\bfseries}
\setbeamertemplate{navigation symbols}{}
\setbeamertemplate{footline}{%
  \leavevmode%
  \hbox{%
    \begin{beamercolorbox}[wd=.5\paperwidth,ht=2.5ex,dp=1.125ex,leftskip=6pt]{author in head/foot}%
      \usebeamerfont{author in head/foot}\textcolor{gray}{\small Úvod do Pythonu na Micro:bitu}
    \end{beamercolorbox}%
    \begin{beamercolorbox}[wd=.5\paperwidth,ht=2.5ex,dp=1.125ex,rightskip=6pt,right]{date in head/foot}%
      \usebeamerfont{date in head/foot}\textcolor{gray}{\small\insertframenumber\,/\,\inserttotalframenumber}
    \end{beamercolorbox}}%
}

% --- Česky korektní uvozovky ---
\usepackage{csquotes}
\newcommand{\uv}[1]{\enquote{#1}}

% --- Zdrojový kód (Python) ---
\usepackage{listings}
\lstset{
  language=Python,
  basicstyle=\ttfamily\small,
  backgroundcolor=\color{code-bg},
  frame=single,
  framerule=0pt,
  xleftmargin=6pt,
  xrightmargin=6pt,
  breaklines=true,
  showstringspaces=false,
  keywordstyle=\color{microbit-blue}\bfseries,
  stringstyle=\color{teal},
  commentstyle=\color{gray}\itshape,
  literate=
    {á}{{\'a}}1 {é}{{\'e}}1 {í}{{\'i}}1 {ó}{{\'o}}1 {ú}{{\'u}}1
    {ů}{{\r{u}}}1 {č}{{\v{c}}}1 {š}{{\v{s}}}1 {ž}{{\v{z}}}1
    {Á}{{\'A}}1 {É}{{\'E}}1 {Í}{{\'I}}1 {Ó}{{\'O}}1 {Ú}{{\'U}}1
    {Č}{{\v{C}}}1 {Š}{{\v{S}}}1 {Ž}{{\v{Z}}}1 {ň}{{\v{n}}}1 {ř}{{\v{r}}}1,
}

% --- Zvýraznění pojmů ---
\newcommand{\pojem}[1]{\textbf{\textcolor{microbit-blue}{#1}}}
\newcommand{\pyinline}[1]{\texttt{#1}}

% --- Grafika a diagramy ---
\usepackage{tikz}

% --- Ostatní ---
\usepackage{graphicx}
\usepackage{hyperref}
\hypersetup{colorlinks=true, urlcolor=microbit-blue}


\documentclass[aspectratio=169]{beamer}

% --- Jazyk a font ---
\usepackage{polyglossia}
\setmainlanguage{czech}
\usepackage{fontspec}

% --- Téma ---
\usetheme{default}
\usecolortheme{default}

\definecolor{microbit-blue}{HTML}{1E4D8C}
\definecolor{microbit-lightblue}{HTML}{D6E4F7}
\definecolor{code-bg}{HTML}{F0F0F0}

\setbeamercolor{frametitle}{bg=microbit-blue, fg=white}
\setbeamercolor{title}{fg=microbit-blue}
\setbeamercolor{block title}{bg=microbit-blue, fg=white}
\setbeamercolor{block body}{bg=microbit-lightblue}
\setbeamercolor{itemize item}{fg=microbit-blue}
\setbeamercolor{itemize subitem}{fg=microbit-blue!70}
\setbeamerfont{frametitle}{size=\large, series=\bfseries}
\setbeamertemplate{navigation symbols}{}
\setbeamertemplate{footline}{%
  \leavevmode%
  \hbox{%
    \begin{beamercolorbox}[wd=.5\paperwidth,ht=2.5ex,dp=1.125ex,leftskip=6pt]{author in head/foot}%
      \usebeamerfont{author in head/foot}\textcolor{gray}{\small Úvod do Pythonu na Micro:bitu}
    \end{beamercolorbox}%
    \begin{beamercolorbox}[wd=.5\paperwidth,ht=2.5ex,dp=1.125ex,rightskip=6pt,right]{date in head/foot}%
      \usebeamerfont{date in head/foot}\textcolor{gray}{\small\insertframenumber\,/\,\inserttotalframenumber}
    \end{beamercolorbox}}%
}

% --- Česky korektní uvozovky ---
\usepackage{csquotes}
\newcommand{\uv}[1]{\enquote{#1}}

% --- Zdrojový kód (Python) ---
\usepackage{listings}
\lstset{
  language=Python,
  basicstyle=\ttfamily\small,
  backgroundcolor=\color{code-bg},
  frame=single,
  framerule=0pt,
  xleftmargin=6pt,
  xrightmargin=6pt,
  breaklines=true,
  showstringspaces=false,
  keywordstyle=\color{microbit-blue}\bfseries,
  stringstyle=\color{teal},
  commentstyle=\color{gray}\itshape,
  literate=
    {á}{{\'a}}1 {é}{{\'e}}1 {í}{{\'i}}1 {ó}{{\'o}}1 {ú}{{\'u}}1
    {ů}{{\r{u}}}1 {č}{{\v{c}}}1 {š}{{\v{s}}}1 {ž}{{\v{z}}}1
    {Á}{{\'A}}1 {É}{{\'E}}1 {Í}{{\'I}}1 {Ó}{{\'O}}1 {Ú}{{\'U}}1
    {Č}{{\v{C}}}1 {Š}{{\v{S}}}1 {Ž}{{\v{Z}}}1 {ň}{{\v{n}}}1 {ř}{{\v{r}}}1,
}

% --- Zvýraznění pojmů ---
\newcommand{\pojem}[1]{\textbf{\textcolor{microbit-blue}{#1}}}
\newcommand{\pyinline}[1]{\texttt{#1}}

% --- Grafika a diagramy ---
\usepackage{tikz}

% --- Ostatní ---
\usepackage{graphicx}
\usepackage{hyperref}
\hypersetup{colorlinks=true, urlcolor=microbit-blue}


\title{Lekce 10: Volný projekt}
\subtitle{Úvod do Pythonu na Micro:bitu}
\date{}

\begin{document}

\begin{frame}
  \titlepage
\end{frame}

% ------------------------------------------------------------------
\begin{frame}{Co všechno umíme}

  Za deset lekcí jsme prošli celý základ Pythonu:

  \bigskip
  \begin{center}
  \begin{tabular}{lll}
    \toprule
    \textbf{Lekce} & \textbf{Python} & \textbf{Micro:bit} \\
    \midrule
    01 & volání funkcí, skript  & \texttt{display.show/scroll} \\
    02 & proměnné, \texttt{int}, \texttt{str} & vizitka \\
    03 & \texttt{if / elif / else} & tlačítka \\
    04 & \texttt{while}, \texttt{for}, \texttt{sleep()} & animace \\
    05 & \texttt{import random} & hod kostkou \\
    06 & \texttt{float}, operátory & senzory \\
    07 & seznamy, \texttt{import music} & melodie \\
    08 & \texttt{def}, parametry, \texttt{return} & refaktoring \\
    09 & řetězce, \texttt{import radio} & komunikace \\
    \bottomrule
  \end{tabular}
  \end{center}
\end{frame}

% ------------------------------------------------------------------
\begin{frame}{Nápady na projekty}

  \begin{columns}[T]
    \begin{column}{0.48\textwidth}
      \begin{block}{Jednodušší (\(\bigstar\))}
        \begin{itemize}
          \item Digitální kostka s animací
          \item Stopky / odpočítávání
          \item Teploměr s varováním
          \item Věštírna s obrázky
        \end{itemize}
      \end{block}
    \end{column}
    \begin{column}{0.48\textwidth}
      \begin{block}{Náročnější (\(\bigstar\bigstar\))}
        \begin{itemize}
          \item Jukebox s vlastními melodiemi
          \item Hádej číslo (micro:bit vs.\ hráč)
          \item Komunikátor přes radio
          \item Hra pro dva hráče (radio)
        \end{itemize}
      \end{block}
    \end{column}
  \end{columns}

  \bigskip
  \begin{exampleblock}{Vlastní nápad?}
    Cokoliv, co tě baví — pokud to funguje a využívá alespoň 3~koncepty z~kurzu.
  \end{exampleblock}
\end{frame}

% ------------------------------------------------------------------
\begin{frame}{Jak na projekt}

  \begin{enumerate}
    \item \textbf{Vyber si téma} — co bude micro:bit dělat?
    \item \textbf{Napiš plán} — jaké koncepty použiješ? Jaké funkce nadefinuješ?
    \item \textbf{Rozděl na části} — začni nejjednodušší částí, postupně rozšiřuj
    \item \textbf{Testuj průběžně} — ověř, že každá část funguje, než přidáš další
    \item \textbf{Odevzdej} — kód + krátký popis co program dělá
  \end{enumerate}

  \bigskip
  \begin{alertblock}{Tip}
    Fungující jednoduchý program je lepší než nefungující složitý.\\
    Začni základem — vylepšovat můžeš vždy.
  \end{alertblock}
\end{frame}

% ------------------------------------------------------------------
\begin{frame}{Co odevzdat}

  \begin{itemize}
    \item[\checkmark] \textbf{Zdrojový kód} (soubor \texttt{.py})
    \item[\checkmark] \textbf{Krátký popis} — co program dělá a které koncepty jsi použil/a
    \item[\checkmark] Volitelně: \textbf{video} s ukázkou programu v~akci
  \end{itemize}

  \bigskip
  Program musí:
  \begin{itemize}
    \item fungovat (být „kompletní", ne jen fragment)
    \item využívat alespoň \textbf{3 koncepty} z~kurzu
    \item obsahovat alespoň \textbf{1 vlastní funkci} (\pyinline{def})
  \end{itemize}

  \bigskip
  \begin{block}{Hodnocení}
    Klíčová otázka: „Fungovalo by to na jiném micro:bitu?"\\
    Pokud ano — splnil/a jsi zadání.
  \end{block}
\end{frame}

\end{document}
